\documentclass[noauthor,noinstructornotes]{ximera}

\title{Week 4}
\author{Math 1193 Design Team}

\begin{document}
\begin{abstract}
    Activities for Week 4.
\end{abstract}
\maketitle

\begin{problem}
Fill in as much of the following table as you can, then 
compare with your group. Are there any cell you cannot fill in?

\begin{center}
\begin{tabular}{|c|c|c|c|} 
 \hline
 \(x\) & \(f(x)\) & \(3f(x)\) & \(f(3x)\) \\
 \hline\hline
 -3 & 0 &  &  \\ 
 \hline
 -1 & -2 &  &  \\
 \hline
 0 & 7 &  &  \\
 \hline
 1 & 3 &  &  \\
 \hline
 3 & -1 &  &  \\
 \hline
\end{tabular}
\end{center}

\begin{instructorNotes}
This is meant to give students a preview of transformations. The rules don't
just come from nowhere. It also is meant to see how students are doing
with function notation. Print out tables and coordinate planes.
\end{instructorNotes}
\end{problem}

\begin{problem}
Fill in as much of the following table as you can, then 
compare with your group. Are there any cells you cannot fill in?

\begin{center}
\begin{tabular}{|c|c|c|c|} 
 \hline
 \(x\) & \(f(x)\) & \(f(x)-1\) & \(f(x - 1)\) \\
 \hline\hline
 1 & 2 &  &  \\ 
 \hline
 2 & 6 &  &  \\
 \hline
 3 & 4 &  &  \\
 \hline
 4 & 9 &  &  \\
 \hline
 5 & -1 &  &  \\
 \hline
\end{tabular}
\end{center}
\end{problem}

\begin{problem}
Fill in as much of the following table as you can, then 
compare with your group. Are there any cells you cannot fill in?

\begin{center}
\begin{tabular}{|c|c|c|c|} 
 \hline
 \(x\) & \(f(x)\) & \(-f(x)\) & \(f(-x)\) \\
 \hline\hline
 -2 & -2 &  &  \\ 
 \hline
 -1 & 0 &  &  \\
 \hline
 0 & 2 &  &  \\
 \hline
 1 & 4 &  &  \\
 \hline
 2 & 6 &  &  \\
 \hline
\end{tabular}
\end{center}
\end{problem}

\begin{problem}
The graph of function \(g\) is a vertical translation of 
the graph of function \(f\).

\begin{center}
    \includegraphics[width=20cm]{transformationsVisual.png}
\end{center}

\begin{enumerate}
    \item Complete the \(g(x)\) column of the table.
    \textbf{Do not fill out the last column yet.}
    \begin{center}
        \begin{tabular}{|c|c|c|c|} 
        \hline
        \(x\) & \(f(x)\) & \(g(x)\) & \(h(x) = f(x) - 2.5\) \\
        \hline\hline
        -4 & 0 &  &  \\ 
        \hline
        -3 & -5.8 &  &  \\
        \hline
        -0.7 & 0 &  &  \\
        \hline
        1.2 & -3.3 &  &  \\
        \hline
        2 & 0 &  &  \\
        \hline
        \end{tabular}
    \end{center}
    \item If \(f(0) = -0.86\), what is \(g(0)\)? Explain to 
    a partner how you know.
    \item Write an equation for \(g(x)\) in terms of \(f(x)\) for any
    input \(x\).
    \item The function \(h\) can be written in terms of \(f\) as 
    \(h(x) = f(x) - 2.5\). Complete the \(h(x)\) column of the 
    table.
\end{enumerate}
\end{problem}

\begin{problem}
In this exercise, we'll draw a graph whose formula is 
difficult to write down concisely.

\begin{enumerate}
    \item Draw the line \(y = 2x - 1\), but only for
    \(-5 < x \le -1\).
    \item On the same coordinate plane, draw the line 
    \(y = -x - 1\), but only for \(-1 < x < 2\).
    \item On the same coordinate plane, draw the line \(y = x\),
    but only for \(2 \le x \le 4\).
    \item Is the resulting graph the graph of a function? Why or
    why not? Discuss with a partner.
\end{enumerate}

\begin{instructorNotes}
Use this as an opportunity to gauge comfort with linear functions.
\end{instructorNotes}
\end{problem}

\end{document}
